\section*{Chapter \ref{ch:b1:comparative-planetology} --- Comparative Planetology}

\begin{solution}{exo:b1:comparative-planetology:1}
$V = \tfrac43\pi(\qty{6.9911e7}{m})^3 = \qty{1.4313e24}{m^3}$ and
$\bar\rho = \num{1898.13e24}/\num{1.4313e24} = \qty{1326}{kg.m^{-3}}$, only
1.33 times the density of water, although Jupiter is 318 times more massive
than the Earth: it is made mostly of the lightest elements.
\end{solution}

\begin{solution}{exo:b1:comparative-planetology:2}
Mercury (radius \qty{2440}{km}, \qty{5429}{kg.m^{-3}}) and Venus
(\qty{6052}{km}, \qty{5243}{kg.m^{-3}}) are small and denser than rock:
\omterm{def:b1:comparative-planetology:planet-types}{terrestrial planets}, of rock and iron. Saturn (\qty{60268}{km},
\qty{687}{kg.m^{-3}}) is huge and lighter than water: a \omterm{def:b1:comparative-planetology:planet-types}{gas giant}, mostly
hydrogen and helium. Uranus (\qty{25559}{km}, \qty{1270}{kg.m^{-3}}) is a
giant but far smaller than Saturn and nearly twice as dense: an \omterm{def:b1:comparative-planetology:planet-types}{ice giant},
whose water, methane and ammonia are heavier than hydrogen.
\end{solution}

\begin{solution}{exo:b1:comparative-planetology:3}
$T_\mathrm{eq} = (2601.3\times0.23/(4\times\num{5.6704e-8}))^{1/4} =
(\num{2.638e9})^{1/4} = \qty{226.6}{K}$. Venus receives 1.9 times more
sunlight than the Earth but its clouds reflect \qty{77}{\percent} of it:
it absorbs only $2601.3\times0.23 = \qty{598}{W.m^{-2}}$ against
$1361\times0.706 = \qty{961}{W.m^{-2}}$ for the Earth.
\end{solution}

\begin{solution}{exo:b1:comparative-planetology:4}
From the most uniform to the most concentrated: the Moon, Mars, the Earth,
Jupiter, Saturn. A point mass at the centre has $I = 0$, hence $C = 0$.
\end{solution}

\begin{solution}{exo:b1:comparative-planetology:5}
$T = 15.945421\times86400 = \qty{1377684}{s}$;
$GM = 4\pi^2(\num{1.22187e9})^3/\num{1377684}^2 = \qty{3.794e16}{m^3.s^{-2}}$;
$M = \num{3.794e16}/\num{6.6743e-11} = \qty{5.685e26}{kg}$, within
\qty{0.03}{\percent} of \qty{568.32e24}{kg}.
\end{solution}

\begin{solution}{exo:b1:comparative-planetology:6}
$S = 1361\times(149.6/228.0)^2 = \qty{585.9}{W.m^{-2}}$ at Mars and
$1361\times(149.6/778.5)^2 = \qty{50.26}{W.m^{-2}}$ at Jupiter, in agreement
with the tables (the small difference at Mars comes from the rounding of the
mean distances).
\end{solution}

\begin{solution}{exo:b1:comparative-planetology:7}
$S \propto d^{-2}$ and $T_\mathrm{eq} \propto S^{1/4}$, so $T_\mathrm{eq}
\propto d^{-1/2}$. For Jupiter, $T_\mathrm{eq} = (50.26\times0.657/(4\times
\num{5.6704e-8}))^{1/4} = \qty{109.8}{K}$. A planet radiating more than it
absorbs must have an internal source of heat: Jupiter would be releasing
heat stored since its formation, or produced as it slowly contracts.
\end{solution}

\begin{solution}{exo:b1:comparative-planetology:8}
Venus: $737 - 226.6 = \qty{510}{K}$, \qty{225}{\percent} of $T_\mathrm{eq}$.
The Earth: $T_\mathrm{eq} = (1361\times0.706/(4\sigma))^{1/4} = \qty{255.1}{K}$,
$288 - 255.1 = \qty{33}{K}$, \qty{13}{\percent}. Mars: $214 - 209.8 =
\qty{4}{K}$, \qty{2}{\percent}.
\end{solution}

\begin{solution}{exo:b1:comparative-planetology:9}
With $\mathrm dm = 4\pi\rho r^2\,\mathrm dr$: $I = \int_0^R\tfrac23r^2\,4\pi
\rho r^2\,\mathrm dr = \tfrac{8\pi}{15}\rho R^5$, and $MR^2 = \tfrac43\pi\rho
R^5$, so $C = \tfrac{8}{15}\cdot\tfrac34 = \tfrac25$. A hollow shell has all
its mass at $r = R$: $I = \tfrac23MR^2$, $C = \tfrac23$.
\end{solution}

\begin{solution}{exo:b1:comparative-planetology:10}
$\tfrac52C\bar\rho = 2.5\times0.35\times5429 = 4750.4$;
$x^2 = (4750.4 - 3300)/(5429 - 3300) = 1450.4/2129 = 0.6812$, $x = 0.8254$,
core radius $0.8254\times2439.7 = \qty{2014}{km}$. $x^3 = 0.5623$ and
$\rho_c = 3300 + 2129/0.5623 = \qty{7086}{kg.m^{-3}}$. The core holds
$0.5623\times7086/5429 = \qty{73}{\percent}$ of the mass. Mercury's core
reaches \qty{83}{\percent} of its radius, against $3485/6371 =
\qty{55}{\percent}$ for the Earth: Mercury is mostly core, with a thin rocky
mantle.
\end{solution}

\begin{solution}{exo:b1:comparative-planetology:11}
(a) The function $t\mapsto t^4$ is convex, so
$\bigl(\tfrac{T_1 + T_2}{2}\bigr)^4 \le \tfrac{T_1^4 + T_2^4}{2} =
T_\mathrm{eq}^4$, with equality only if $T_1 = T_2$. (b) $T_1^4 =
2\times270^4 - 100^4 = \num{1.0529e10}$, $T_1 = \qty{320}{K}$, and the mean
temperature is $(320 + 100)/2 = \qty{210}{K}$, \qty{60}{K} below
$T_\mathrm{eq}$.
\end{solution}

\begin{solution}{exo:b1:comparative-planetology:12}
$m = 4\pi R^2p/g = 4\pi\times(\num{6.0518e6})^2\times\num{9.2e6}/8.9 =
\num{4.602e14}\times\num{1.034e6} = \qty{4.76e20}{kg}$: 93 times the mass of
the Earth's atmosphere, and \num{9.8e-5} of the mass of Venus. The Earth
holds a comparable amount of carbon, not in its air but in carbonate rocks
(limestones) and in the ocean, where liquid water allowed it to be stored.
\end{solution}

\begin{solution}{pb:b1:comparative-planetology:1}
\textbf{1.} $T = 0.31891\times86400 = \qty{27553.8}{s}$.

\textbf{2.} $GM = 4\pi^2\times(\num{9.378e6})^3/\num{27553.8}^2 =
\qty{4.289e13}{m^3.s^{-2}}$.

\textbf{3.} $M = \num{4.289e13}/\num{6.6743e-11} = \qty{6.426e23}{kg}$,
within \qty{0.14}{\percent} of the tabulated \qty{6.4169e23}{kg}.

\textbf{4.} $\num{1.06e16}/\num{6.426e23} = \num{1.6e-8}$: Phobos is
negligible.

\textbf{5.} $v = 2\pi a/T = 2\pi\times\num{9.378e6}/27553.8 =
\qty{2138}{m.s^{-1}}$, about \qty{2.1}{km.s^{-1}}.

\textbf{6.} $T = \qty{7.65}{h}$, 3.2 times shorter than the rotation of
Mars: Phobos goes round faster than the ground turns beneath it, so it
overtakes the rotation. Seen from Mars, it rises in the west and sets in
the east.

\textbf{7.} $V = \tfrac43\pi(\num{3.3895e6})^3 = \qty{1.6312e20}{m^3}$.

\textbf{8.} $\bar\rho = \num{6.426e23}/\num{1.6312e20} = 3939$, that is
\qty{3940}{kg.m^{-3}}.

\textbf{9.} $3940/5514 = 0.715$. The Earth, larger and more massive,
compresses its own interior more strongly; and its iron core is a larger
share of the planet.

\textbf{10.} $g = GM/R^2 = \num{4.289e13}/(\num{3.3895e6})^2 =
\qty{3.73}{m.s^{-2}}$, \qty{38}{\percent} of the Earth's.

\textbf{11.} $v_\mathrm{esc} = \sqrt{2GM/R} = \sqrt{2\times\num{4.289e13}/
\num{3.3895e6}} = \qty{5.03}{km.s^{-1}}$.

\textbf{12.} $T_\mathrm{eq} = (586.2\times0.75/(4\times\num{5.6704e-8}))^{1/4}
= \qty{209.8}{K}$.

\textbf{13.} $214 - 209.8 = \qty{4}{K}$: a weak greenhouse warming.

\textbf{14.} $m = p/g = 636/3.733 = \qty{170}{kg.m^{-2}}$; for the Earth
$101400/9.82 = \qty{10326}{kg.m^{-2}}$, 61 times more.

\textbf{15.} $4\pi R^2\times170.4 = \num{1.444e14}\times170.4 =
\qty{2.46e16}{kg}$.

\textbf{16.} Above each square metre there is sixty times less gas than on
the Earth, so the infrared absorption, although by carbon dioxide, is weak;
and the air of Mars holds very little water vapour, the other major
absorber.

\textbf{17.} A uniform Mars would have $C = 0.4$; the measured 0.366 shows
that its density increases toward the centre: Mars has a dense core.

\textbf{18.} $\bar\rho = \rho_m + (\rho_c - \rho_m)x^3$.

\textbf{19.} $C = \tfrac25[\rho_m + (\rho_c - \rho_m)x^5]/\bar\rho$.

\textbf{20.} From 18 and 19, $(\rho_c - \rho_m)x^3 = \bar\rho - \rho_m$ and
$(\rho_c - \rho_m)x^5 = \tfrac52C\bar\rho - \rho_m$; dividing gives $x^2$.

\textbf{21.} $\tfrac52C\bar\rho = 2.5\times0.366\times3940 = 3605.1$;
$x^2 = (3605.1 - 3500)/(3940 - 3500) = 105.1/440 = 0.2389$, $x = 0.4887$.
$x^3 = 0.1167$ and $\rho_c = 3500 + 440/0.1167 = \qty{7269}{kg.m^{-3}}$.

\textbf{22.} $x^3\rho_c/\bar\rho = 0.1167\times7269/3940 = 0.215$, that is
\qty{21.5}{\percent} of the mass.

\textbf{23.} With $\rho_m = 3450$: $x^2 = 155.1/490 = 0.3165$, $x = 0.5626$,
core radius \qty{1907}{km}, $\rho_c = \qty{6202}{kg.m^{-3}}$. With $\rho_m =
3550$: $x^2 = 55.1/390 = 0.1413$, $x = 0.3759$, \qty{1274}{km},
$\rho_c = \qty{10894}{kg.m^{-3}}$. A change of \qty{1.4}{\percent} in the
assumed mantle density moves the core radius by several hundred kilometres:
the mean density and $C$ alone cannot fix the core of Mars precisely, which
is why seismometers were sent there.

\textbf{24.} $0.4887\times3389.5 = \qty{1657}{km}$: the radius of the core
of Mars from the model is \textbf{about \qty[text-series-to-math=true]{1660}{km}}, within the recent
seismic range of \qty{1650}{km} to \qty{1675}{km} and well below the
\qty{1830}{km} of the planetary table --- an agreement that question 23
shows to be partly luck in the choice of $\rho_m$.
\end{solution}
